\documentclass[11pt]{article}

\usepackage[a4paper,margin=1in]{geometry}
\usepackage[T1]{fontenc}
\usepackage[utf8]{inputenc}
\usepackage{lmodern}
\usepackage{amsmath,amssymb}
\usepackage{physics}
\usepackage{hyperref}
\usepackage{enumitem}
\usepackage{setspace}

\hypersetup{colorlinks=true,linkcolor=blue,citecolor=blue,urlcolor=blue}
\setstretch{1.1}

\title{Saturated Superfluid Vacuum VI-a\\
Galactic Standing Waves and Flat Rotation Curves}
\author{Stig Norland}
\date{Draft}

\begin{document}
\maketitle

\begin{abstract}
This paper studies flat rotation curves in the Saturated Superfluid Vacuum (SSV) framework as
a consequence of galactic standing-wave structure. The central claim is that galactic discs are
coherent standing-wave or solitonic structures anchored to the central black-hole condensate,
and that this structure produces flat rotation curves without dark-matter halos.

The paper is intentionally narrow. It focuses on the standing-wave scale, the disc structure,
and the rotation-curve mechanism, together with a concrete application to a selected galaxy.
Broader cosmological claims are left aside.
\end{abstract}

\section{Introduction}

This paper studies flat rotation curves in SSV as a consequence of galactic standing-wave
structure.

\section{Central Source Scale}

Introduce the central black-hole scale and its role in setting the galactic standing-wave scale.

\section{Disc Soliton and Standing Wave}

Present the disc as a coherent structure of the medium.

\section{Rotation-Curve Mechanism}

Explain how the radial structure produces flat or nearly flat rotation curves.

\section{Application to Andromeda}

Provide a focused worked example.

\section{Predictions and Limits}

Keep predictions local to the mechanism developed in this paper.

\section{Conclusion}

The flat-curve claim should stand or fall here as a focused galactic-structure argument, not
as part of a broader cosmological manifesto.

\end{document}
